\section{Galois representations and local parameters}\label{sec:galois}

\providecommand{\galoisWD}{\mathrm{WD}}
\providecommand{\galoisIrr}{\operatorname{Irr}}
\providecommand{\galoisrc}{\mathrm{rc}}
\providecommand{\galoisSp}{\operatorname{Sp}}
\providecommand{\galoisHT}{\mathrm{HT}}
\providecommand{\galoisinf}{\operatorname{inf}}
\providecommand{\galoisrecT}{\mathrm{recT}}
\providecommand{\galoisU}{\mathrm{U}}

This section is devoted to the two inputs of the proof that involve $\ell$-adic Galois
representations and local Langlands parameters:
\begin{itemize}[nosep]
\item $\sfname{RecTExists}$ (\cref{def:rectexists}, \cref{thm:rectexists}): at every finite place
  $w$ of every number field $E$ there is a \emph{local parameter map}, i.e.\ a map from the
  irreducible smooth representations of the groups $\GL_n(E_w)$, $n\ge0$, to Frobenius-semisimple
  Weil--Deligne representations which obeys the twelve laws (L1)--(L12) of \cref{def:rectexists};
\item $\sfname{Row0b}(\galoisrecT)$ (\cref{def:row0b}, \cref{thm:row0b}): every regular algebraic,
  polarized, cuspidal automorphic representation $\pi$ of $\GL_a(\A_E)$, $E$ a CM field, has an
  $\ell$-adic Galois representation with the expected Hodge--Tate weights which is compatible with
  $\pi$, in a partial sense, with respect to \emph{every} local parameter map.
\end{itemize}
These are the only inputs of the top-level deduction (\cref{sec:top}) involving Galois
representations and local Langlands parameters. \Cref{thm:row0b} is the input from which the window
bound $\sfname{VWin}$ is derived (\cref{sec:window}), and local parameter maps are used in the
construction of the hypothesis package (\cref{sec:families}).

Both statements are modelled on published theorems, the local Langlands correspondence for $\GL_n$
\cite{HarrisTaylor2001,Henniart2000,Scholze2013} and the existence of Galois representations attached
to regular algebraic polarized cuspidal representations with local--global compatibility
\cite{HarrisTaylor2001,TaylorYoshida2007,Shin2011,ChenevierHarris2013,Caraiani2012,Caraiani2014,BarnetLambEtAl2014},
and both are weaker than these, as \cref{rem:galois:llc,rem:galois:blggt} make precise. The formal
proof \cite{LeanSource2026} uses neither published theorem: local parameter maps are
constructed explicitly (\cref{sec:galois:llcproof}), and the Galois representations are constructed
(\cref{sec:galois:arch}) through totally definite unitary groups, by a geometric realization in the
cohomology of Kuga--Sato varieties over Shimura varieties of Harris--Taylor type (\cref{thm:geoone}),
$\ell$-adic congruences and limits, and theta lifts between definite unitary groups
(\cref{thm:galois:tlift}), followed by descent and base change to $\GL_a$.

% ---------------------------------------------------------------------------
\subsection{Local parameter maps}\label{sec:galois:local}

In \cref{sec:galois:local,sec:galois:llcproof}, $E$ is a number field, $\Omega$ an algebraic closure
of $E$ and $G_E=\Aut_E(\Omega)$ with its Krull topology. A finite place of $E$ is a nonzero prime $w$
of $\OO_E$; we write $F=E_w$ for the completion, $\OO_F$, $k_F$, $\ord_F$ and $|\cdot|$ for its
valuation ring, residue field, normalized valuation and normalized absolute value, and $q=\Nm w=|k_F|$.
A \emph{place of $\Omega$ above $w$} is a valuation ring $A\subset\Omega$ with
$A\cap E=\OO_{E,(w)}$. It determines decomposition, inertia and wild inertia groups
$D_A\supset I_A\supset P_A$ in $G_E$, and the \emph{Weil group} $W_A\subset D_A$ of the elements
$\sigma$ that act on the residue field of $A$ as $x\mapsto x^{q^{\deg\sigma}}$ for some
$\deg\sigma\in\Z$. As $D_A\cong\Gal(\overline F/F)$, these are models of the Galois, inertia and Weil
groups of $F$; arithmetic Frobenius elements have degree $1$, and $I_A$ is the kernel of $\deg$.

\begin{definition}[Weil--Deligne representations, local Artin maps]\label{def:galois:wd}
\begin{enumerate}[label=(\alph*),nosep]
\item $\galoisWD_A$ is the set of isomorphism classes of Frobenius-semisimple Weil--Deligne
  representations $(\rho,N)$ on the spaces $\C^n$, $n\ge0$: $\rho\colon W_A\to\GL_n(\C)$ is a
  homomorphism, trivial on an open subgroup of $I_A$, with $\rho(\sigma)$ semisimple for all
  $\sigma$, and $N$ is nilpotent with $\rho(\sigma)N\rho(\sigma)^{-1}=q^{\deg\sigma}N$. For $s\in\R$ we
  write $V(s)=V\otimes(\sigma\mapsto q^{s\deg\sigma})$; $V^\vee$ is the dual, $\langle\xi\rangle$ the
  one-dimensional representation of a smooth character $\xi$ of $W_A$, and $\galoisSp(m)$ is $\C^m$
  with $\rho(\sigma)=\mathrm{diag}\bigl(q^{(i-(m-1)/2)\deg\sigma}\bigr)_{0\le i<m}$ and $Ne_i=e_{i+1}$.
  $V$ is \emph{tame} if $\rho(P_A)=1$, and \emph{abelian} if $\rho$ factors through $W_A^{\mathrm{ab}}$.
\item A \emph{local Artin map} is a surjective homomorphism $a\colon W_A\to F^\times$, continuous on
  $I_A$, with $\ord_F(a(\sigma))=-\deg\sigma$, such that every smooth character of $W_A$ is $\chi\circ
  a$ for a smooth character $\chi$ of $F^\times$, and such that $\mathrm{art}(a(\sigma))=\sigma^{\mathrm{ab}}$
  for every global Artin map $\mathrm{art}\colon\A_E^\times\to G_E^{\mathrm{ab}}$ (with $F^\times\subset
  \A_E^\times$ at $w$). Thus $a$ is the inverse of the local reciprocity map, normalized so that
  uniformizers correspond to geometric Frobenius elements, as in \cite{Deligne1973,HarrisTaylor2001}.
\end{enumerate}
\end{definition}

\begin{definition}[Representations of $\GL_n(F)$; good representations]\label{def:galois:irr}
$\galoisIrr(F)$ is the disjoint union over $n\ge0$ of the sets of isomorphism classes of smooth
irreducible complex representations of $\GL_n(F)$. For $x\in\galoisIrr(F)$ we write $n(x)$ for its
index, $x\nu^s=x\otimes|\det|^s$, $x\otimes\chi=x\otimes(\chi\circ\det)$ for smooth characters $\chi$
of $F^\times$, $\omega_x$ for its central character and $x^\vee$ for its contragredient; \emph{cuspidal}
means supercuspidal, and $x$ has \emph{depth zero} if it has a nonzero vector fixed by
$K_1=\ker(\GL_n(\OO_F)\to\GL_n(k_F))$. For a cuspidal $\rho$ and $m\ge1$, $\Delta=[\rho,\rho\nu^{m-1}]$
is a \emph{segment}, and for a finite multiset $\mathfrak m$ of segments $L(\mathfrak m)\in\galoisIrr(F)$
denotes its Langlands quotient in the form of Zelevinsky and Rodier \cite{Zelevinsky1980,Rodier1982}:
the irreducible cosocle of the product, taken in a standard order, of the essentially
square-integrable representations $L(\Delta)$, $\Delta\in\mathfrak m$. Every $x$ is $L(\mathfrak m)$ for a
unique $\mathfrak m$, and the supercuspidal support of $L(\mathfrak m)$ is the multiset of the terms
$\rho,\dots,\rho\nu^{m-1}$ of its segments.

A cuspidal $y\in\galoisIrr(F)$ is \emph{good} if $n(y)=1$, or if $F\cong\Q_g$ as topological fields
for some prime $g$ and $y=y_0\otimes\chi$ with $y_0$ of depth zero and $\chi$ smooth. An arbitrary
$x\in\galoisIrr(F)$ is \emph{good} if its supercuspidal support consists of good cuspidal
representations.
\end{definition}

Thus, if $[F:\Q_p]>1$, the good representations are the irreducible subquotients of principal
series; if $F\cong\Q_p$, the supercuspidal support of a good representation may also contain twists,
by arbitrary (possibly wildly ramified) smooth characters, of depth-zero supercuspidal representations.

\begin{definition}[Local parameter maps; $\sfname{RecTExists}$]\label{def:rectexists}
Let $a$ be a local Artin map. A map $\galoisrc\colon\galoisIrr(F)\to\galoisWD_A$ satisfies
$\sfname{LawsT}(a,\galoisrc)$ if the following laws hold for all $x,x',y\in\galoisIrr(F)$, all smooth
characters $\chi$ of $F^\times$ and all $s\in\R$:
\begin{enumerate}[label=(L\arabic*),leftmargin=3.2em,nosep]
\item $\dim\galoisrc(x)=n(x)$;
\item $\galoisrc(\chi)=\langle\chi\circ a\rangle$, where $\chi$ is viewed as an element of
  $\galoisIrr(F)$ with $n=1$;
\item $\galoisrc\bigl(L(\Delta_1,\dots,\Delta_k)\bigr)=\bigoplus_{i=1}^k\galoisrc(\rho_i)\bigl(\tfrac{m_i-1}2\bigr)
  \otimes\galoisSp(m_i)$ for segments $\Delta_i=[\rho_i,\rho_i\nu^{m_i-1}]$;
\item $\galoisrc(x)$ is irreducible if $x$ is cuspidal;
\item $\galoisrc(x\nu^s)=\galoisrc(x)(s)$;
\item $|\det\galoisrc(x)(\sigma)|=|\omega_x(a(\sigma))|$ for all $\sigma\in W_A$;
\end{enumerate}
if moreover $x$ and $y$ are good:
\begin{enumerate}[label=(L\arabic*),leftmargin=3.2em,start=7,nosep]
\item $\galoisrc(x\otimes\chi)=\galoisrc(x)\otimes\langle\chi\circ a\rangle$;
\item $\det\galoisrc(x)=\omega_x\circ a$;
\item $\galoisrc(x^\vee)=\galoisrc(x)^\vee$;
\item if $\galoisrc(x')=\galoisrc(y)$ then $x'=y$;
\item $\galoisrc(x)$ is tame if $x$ has depth zero;
\end{enumerate}
and the \emph{type law}:
\begin{enumerate}[label=(L\arabic*),leftmargin=3.2em,start=12,nosep]
\item suppose that $F\cong\Q_g$, $n\ge1$, $\theta$ is a regular character of $\F_{g^n}^\times$ (that
  is, $\theta^{g^i}\ne\theta$ for $0<i<n$), $\tau$ is a cuspidal representation of $\GL_n(\F_g)$ whose
  character is given by Green's formula $\tr\tau(y)=(-1)^{n-1}\sum_{j<n}\theta(z^{g^j})$ at every $y$
  with irreducible characteristic polynomial and every root $z\in\F_{g^n}^\times$ of it, and
  $\sigma\in\galoisIrr(F)$ with $n(\sigma)=n$ contains the inflation of $\tau$ to $\GL_n(\Z_g)$. Then
  $\galoisrc(\sigma)$ has $N=0$, and for every $t\in I_A$ the characteristic polynomial of
  $\galoisrc(\sigma)(t)$ is $\prod_{j<n}\bigl(X-\theta(z_t^{g^j})\bigr)$, where $z_t\in\F_{g^n}^\times$ is
  the residue of $a_n(t)$ for the local Artin map $a_n$ of the unramified extension of $F$ of
  degree $n$.
\end{enumerate}
A map $\galoisrc$ is a \emph{local parameter map} ($\galoisrecT(\galoisrc)$; $\sfname{IsRecT}$ in the
development) if $\sfname{LawsT}(a,\galoisrc)$ holds for some local Artin map $a$. $\sfname{RecTExists}$
is the proposition: for every number field $E$, algebraic closure $\Omega$, finite place $w$ of $E$
and place $A$ of $\Omega$ above $w$, there is a local parameter map $\galoisIrr(E_w)\to\galoisWD_A$.
\end{definition}

In the formal statement of (L12) the unramified extension is realized as the completion of an
auxiliary number field $E\subset E'\subset\Omega$ at the place below $A$; the law is required for all
such realizations.

\begin{theorem}[Existence of local parameter maps]\Fmark\label{thm:rectexists}
$\sfname{RecTExists}$ holds. \formal{UpperTheory.Zero.LLC.recTExists}
\end{theorem}
%P2PTR thm:rectexists
\begin{proof}See \cref{sec:pf:llc}.
\end{proof}


\begin{proposition}[Uniqueness on good representations]\Fmark\label{prop:galois:unique}
Any two local parameter maps $\galoisrc,\galoisrc'\colon\galoisIrr(E_w)\to\galoisWD_A$ agree on all
good representations. \formal{UpperTheory.Zero.LLC.recTUnique}
\end{proposition}

\begin{proofoutline}
The local Artin map is unique (Step~0 of \cref{sec:galois:llcproof}), so both maps obey the laws for
the same $a$; on characters both are given by (L2). For $F\cong\Q_g$ and $y$ cuspidal of depth zero,
(L1), (L4), (L8), (L11) and (L12) say that $\galoisrc(y)$ and $\galoisrc'(y)$ are irreducible, tame, of
dimension $n(y)$, with determinant $\omega_y\circ a$ and the restriction to $I_A$ prescribed by (L12),
which determines them (Step~2 of \cref{sec:galois:llcproof}). Then (L7) and (L3) give the claim.

%P2PTR prop:galois:unique
A complete proof is given in \cref{sec:pf:llc}.
\end{proofoutline}

\begin{remark}[Comparison with the local Langlands correspondence]\label{rem:galois:llc}
The local Langlands correspondence $\mathrm{rec}_F$ \cite{HarrisTaylor2001,Henniart2000,Scholze2013} is
a bijection from the irreducible smooth representations of $\GL_n(F)$ onto the $n$-dimensional
elements of $\galoisWD_A$, characterized \cite{Henniart1993} by class field theory, compatibility with
twists, central characters and contragredients, and the preservation of $L$- and $\varepsilon$-factors
of pairs; it satisfies (L3). The laws are much weaker: (i) no surjectivity is required, and
injectivity only in the form (L10); (ii) no condition on $L$- or $\varepsilon$-factors of pairs is
imposed, although for $\mathrm{rec}_F$ these determine the parameters of supercuspidal
representations; (iii) arbitrary twists, central characters, contragredients and the implication
``depth zero $\Rightarrow$ tame'' are controlled only on good representations, while for arbitrary $x$
only the twists $\nu^s$ (L5) and $|\omega_x|$ (L6) are; (iv) cuspidal representations go to irreducible
ones, but not conversely. On
the other hand (L12) fixes, through Green's parametrization, the restriction to inertia of the
parameters of depth-zero supercuspidal representations of $\GL_n(\Q_g)$. Thus a local parameter map is
determined on good representations (\cref{prop:galois:unique}), where it is given by explicit formulas
(on toral representations, like $\mathrm{rec}_F$, by class field theory and (L3)). On a cuspidal $x$ that
is not good the laws only ask that $\galoisrc(x)$ be irreducible of dimension $n(x)$ with
$|\det\galoisrc(x)|=|\omega_x\circ a|$, compatibly with the twists $\nu^s$, and distinct from the values
at good representations; the map constructed below gives all such $x$ with the same $n(x)$ and
$|\omega_x|$ one fixed auxiliary value (Step~4), so it is not $\mathrm{rec}_F$ there. (Maps satisfying
the development's formalization of the full correspondence satisfy (L1)--(L11); this is not used.)
\end{remark}

\begin{remark}[Why the weak laws suffice]\label{rem:galois:suffices}
Local parameter maps enter the proof only through the laws and \cref{prop:galois:unique}. In
$\sfname{Row0b}(\galoisrecT)$ compatibility is required for every local parameter map, but
$\sfname{AwayB}$ (\cref{def:galois:awayb}) constrains $y=\galoisrc(\sigma^\natural)$, $\sigma^\natural=
\sigma\otimes|\det|^{(1-a)/2}$ for a local component $\sigma$, only if $N_y=0$ and $y$ is either
abelian, or tame and squarefree at a place with $E_w\cong\Q_g$. In the first case (L1), (L3), (L4)
force the support of $\sigma^\natural$ to consist of characters; in the second, the proof of
\cref{thm:row0b} shows, using (L11), (L12) and a count of depth-zero supercuspidal representations,
that $\sigma$ has depth zero. Either way $\sigma^\natural$ is good, and the requirement does not depend
on $\galoisrc$. The auxiliary values of Step~4 below are irreducible of dimension $\ge2$ and wildly
ramified, so for them only $\dim x=\dim y$ is required. Elsewhere (notably in \cref{sec:families}) a
local parameter map is chosen at each place, and only its laws are used.
\end{remark}

% ---------------------------------------------------------------------------
\subsection{Construction of local parameter maps}\label{sec:galois:llcproof}

The proof of \cref{thm:rectexists} is an explicit construction; no result on the local Langlands
correspondence itself is used.

\begin{proofoutline}
Fix $E,\Omega,w,A$, and put $F=E_w$ and $q=\Nm w$.

\emph{Step 0: local class field theory.} There is exactly one local Artin map $a$. Existence is local
class field theory \cite{Serre1967,ArtinTate1968,Tate1979}, which the development derives from global
Artin reciprocity. For uniqueness, two local Artin maps have the same composite with
$F^\times\to\A_E^\times\to G_E^{\mathrm{ab}}$, which is injective on $F^\times$ by Chevalley's theorem on
congruence subgroups of units \cite{Chevalley1951}.

\emph{Step 1: characters.} Put $r(\chi)=\langle\chi\circ a\rangle$.

\emph{Step 2: depth-zero cuspidal representations.} Assume $F\cong\Q_g$ and let $y$ be cuspidal of
depth zero, $n=n(y)$. The $K_1$-invariants of $y$ form a cuspidal irreducible representation $\tau$ of
$\GL_n(\F_g)$, unique up to isomorphism, and $y$ is determined by $\tau$ and $\omega_y$
\cite{MoyPrasad1996,BushnellKutzko1993}. By Green's theorem \cite{Green1955}, $\tau$ satisfies Green's
formula for a regular character $\theta$ of $\F_{g^n}^\times$, unique up to $\theta\mapsto\theta^{g^j}$
(the development deduces the uniqueness from an elementary lemma on the sums
$\sum_{j<n}\theta(z^{g^j})$). Reducing the local Artin map of the unramified extension of degree $n$
gives a surjective character $R\colon I_A\to\F_{g^n}^\times$, trivial on $P_A$, with
$R(\varphi t\varphi^{-1})=R(t)^g$ for $\deg\varphi=1$ and $\Nm_{\F_{g^n}/\F_g}R(t)\equiv a(t)\bmod g$. Let
$r(y)$ be the representation induced from the Weil group of that extension: on the functions
$f\colon\Z\to\C$ with $f(j+n)=b\,f(j)$, $t\in I_A$ acts by $(tf)(j)=\theta^{g^j}(R(t))f(j)$ and a fixed
$\varphi$ of degree $1$ by $(\varphi f)(j)=f(j+1)$, with $b$ such that $\det r(y)(\varphi)=\omega_y(a(\varphi))$.
It is irreducible ($\theta$ is regular), tame, with $N=0$ and
determinant $\omega_y\circ a$ (on $I_A$ because scalars $u\in\Z_g^\times$ act on $\tau$ by
$\theta(u\bmod g)$), and it satisfies (L12). Since an irreducible tame representation is determined by
its restriction to inertia and its determinant, $r(y)$ is characterized by these properties, and
$y\mapsto r(y)$ is injective on depth-zero cuspidal representations.

\emph{Step 3: twists.} For a good cuspidal $y=y_0\otimes\chi$ with $y_0$ of depth zero put
$r(y)=r(y_0)\otimes\langle\chi\circ a\rangle$. This is well defined: if $y_0\otimes\chi=y_1\otimes\chi_1$
then $\chi\chi_1^{-1}$ is tame, and $r(y_0\otimes\chi')=r(y_0)\otimes\langle\chi'\circ a\rangle$ for tame
$\chi'$, because twisting $\tau$ by $\bar\chi'\circ\det$ multiplies $\theta$ by
$\bar\chi'\circ\Nm_{\F_{g^n}/\F_g}$.

\emph{Step 4: auxiliary values.} For each $n\ge2$ fix an irreducible $J_n\in\galoisWD_A$ of dimension
$n$ with $|\det J_n|=1$ on which $P_A$ does not act by scalars, so that $J_n$ is not a twist of a tame
representation (it is induced, as in Step~2, from wildly ramified characters of $I_A$ built from an
additive character of the residue field of the unramified extension of degree $n$). For every
cuspidal $x$ which is not good (so $n(x)\ge2$) put $r(x)=J_{n(x)}(e(x))$, where
$|\omega_x(z)|=|z|^{n(x)e(x)}$.

\emph{Step 5: extension along the Langlands classification.} Define $\galoisrc$ by the formula (L3)
with $r$ on the right-hand side; this is well defined by the uniqueness of $\mathfrak m$, and
$\galoisrc=r$ on cuspidal representations. The laws hold for cuspidal representations by Steps 1--4
((L10) for the values of Step~4 because $J_n$ is not a twist of a tame representation), and they
propagate: (L1), (L6), (L8) because the support of $L(\mathfrak m)$ consists of the terms of its segments
and the central character is the product of those of the support; (L5), (L7) because twisting
commutes with Langlands data; (L9) because $L(\mathfrak m)^\vee=L(\mathfrak m^\vee)$ with
$[\rho,\rho\nu^{m-1}]^\vee=[(\rho\nu^{m-1})^\vee,\rho^\vee]$; (L10) because a decomposition
$\bigoplus_iV_i\otimes\galoisSp(m_i)$ with irreducible $V_i$ is unique up to order; (L11) because the
support of a depth-zero representation consists of depth-zero representations \cite{MoyPrasad1996};
and (L12) because a representation containing an inflated cuspidal type is cuspidal of depth zero.

The classical inputs (global Artin reciprocity, Chevalley's theorem, the results of Bernstein and
Zelevinsky \cite{BernsteinZelevinsky1976,BernsteinZelevinsky1977}, the Langlands classification,
Zelevinsky's description of the Grothendieck ring and the contragredients of Langlands data
\cite{Zelevinsky1980}, the depth-zero support theorem \cite{MoyPrasad1996} and Green's theorem
\cite{Green1955}) are all proved in the formal development.
\end{proofoutline}

% ---------------------------------------------------------------------------
\subsection{Galois representations attached to polarized cuspidal representations}\label{sec:galois:row}

From now on $E$ is a CM field with maximal totally real subfield $E^+$ and complex conjugation $c$,
$\ell$ is a prime and $\iota\colon\Qlbar\to\C$ is a field isomorphism.

\begin{definition}[Automorphic notions]\label{def:galois:aut}
Let $a\ge1$.
\begin{enumerate}[label=(\alph*),nosep]
\item A \emph{cuspidal automorphic representation} of $\GL_a(\A_E)$ is an irreducible
  $(\mathfrak g_\infty,K_\infty)\times\GL_a(\A_{E,f})$-submodule $V_\pi$ of the space of smooth cusp forms
  on $\GL_a(E)\backslash\GL_a(\A_E)$ (of moderate growth, $K_\infty$-finite, and finite under the
  centre of the enveloping algebra, which is modelled by Capelli elements); no unitarity or central
  character condition is imposed. A \emph{local component} of $\pi$ at a finite place $w$ is a smooth
  irreducible representation $\sigma$ of $\GL_a(E_w)$ with a nonzero equivariant map $\sigma\to V_\pi$;
  we put $\sigma^\natural=\sigma\otimes|\det|^{(1-a)/2}\in\galoisIrr(E_w)$.
\item For $\tau\colon E\to\C$, $\galoisinf_\tau(\pi)$ is the multiset of $a$ complex numbers $z$ such
  that the Capelli element $\mathrm{cap}_\tau(u)$ acts on $V_\pi$ by $\prod_z(u+z)$ for all $u\in\C$; it
  is the Harish-Chandra parameter of $\pi_\infty$ at $\tau$, normalized so that the trivial
  representation has parameter $\{\frac{a-1}2,\frac{a-3}2,\dots,\frac{1-a}2\}$. $\pi$ is \emph{regular
  algebraic} if for every $\tau$ the multiset $\galoisinf_\tau(\pi)$ consists of $a$ distinct elements
  of $\frac{a-1}2+\Z$ \cite{Clozel1990}.
\item For a Hecke character $\chi\colon\A_{E^+}^\times/(E^+)^\times\to\C^\times$, the pair $(\pi,\chi)$ is
  \emph{polarized} if $\chi_v(-1)=(-1)^a$ at every infinite place $v$ of $E^+$ and
  $\pi^c\cong\pi^\vee\otimes(\chi\circ\Nm_{E/E^+}\circ\det)$, the latter being expressed by a
  nondegenerate invariant pairing. This is the notion of \cite[\S2.1]{BarnetLambEtAl2014}.
\end{enumerate}
\end{definition}

\begin{definition}[Galois-side notions]\label{def:galois:awayb}
Let $r\colon G_E\to\GL_a(\Qlbar)$ be a homomorphism, $w$ a finite place of $E$ and $A$ a place of
$\Omega$ above $w$.
\begin{enumerate}[label=(\alph*),nosep]
\item If $w\nmid\ell$, $\galoisWD_\iota(r|_{W_A})\in\galoisWD_A$ is the Frobenius-semisimplified
  Weil--Deligne representation attached to $r|_{W_A}$ by Grothendieck's $\ell$-adic monodromy theorem,
  transported to $\C$ by $\iota$ (with a default value if there is none).
\item $y\in\galoisWD_A$ is \emph{squarefree} if, for every $\gamma\in I_A$ which together with $P_A$
  topologically generates $I_A$, every multiple root of the characteristic polynomial of $y(\gamma)$
  is a $(\Nm w-1)$-th root of unity. The place $w$ is \emph{$n$-large} if $E_w\cong\Q_g$ as topological
  fields for a prime $g$ such that $n=1$ or $g>16n^2$.
\item For $x,y\in\galoisWD_A$, $\sfname{AwayB}(x,y)$ means: $\dim x=\dim y$; if $y$ is abelian and
  $N_y=0$, then $x=y$; and if $w$ is $(\dim y)$-large and $y$ is tame, squarefree and has $N_y=0$,
  then $N_x=0$ and $x$ and $y$ have the same characteristic polynomials on $I_A$.
\item If $w\mid\ell$ we use the usual notions for representations of $D_A$ \cite{Fontaine1994}:
  crystalline, de Rham, potentially semistable, the $\tau$-Hodge--Tate weights (jumps of the Hodge
  filtration along $\tau\colon E\to\Qlbar$ inducing $w$), and Fontaine's Weil--Deligne representation
  $\galoisWD(\rho)$, transported to $\C$ by $\iota$. (In the development they are expressed relative to
  any package of period rings identified with its $B_{\mathrm{cris}}$ and $B_{\mathrm{dR}}$, and any
  identification of $D_A$ with $\Gal(\Qlbar/K)$, $K/\Ql$ finite with $E$ dense in $K$; the clauses
  below are required for all such choices.)
\end{enumerate}
\end{definition}

\begin{definition}[$\sfname{PropsB}$ and $\sfname{Row0b}$]\label{def:row0b}
A \emph{recognition notion} $\mathrm{Rc}$ is a predicate on the maps
$\galoisrc\colon\galoisIrr(E_w)\to\galoisWD_A$ (for all number fields $E$ and all $w$, $A$); the
instance used is $\galoisrecT$ (\cref{def:rectexists}). Let $\pi$ be a cuspidal automorphic
representation of $\GL_a(\A_E)$, $r\colon G_E\to\GL_a(\Qlbar)$ a homomorphism, $\mathrm{wt}\in\Z$, and
$\galoisHT=(\galoisHT_\tau)_\tau$ a family of multisets of integers indexed by the embeddings
$\tau\colon E\to\Qlbar$. Then $\sfname{PropsB}_{\mathrm{Rc}}(\pi,\iota,r,\mathrm{wt},\galoisHT)$ is the
conjunction of:
\begin{description}[nosep]
\item[(cont, ss)] $r$ is continuous and semisimple;
\item[(away)] for every finite place $w\nmid\ell$, every place $A$ above $w$, every $\galoisrc$ with
  $\mathrm{Rc}(\galoisrc)$ and every local component $\sigma$ of $\pi$ at $w$:
  $\sfname{AwayB}\bigl(\galoisWD_\iota(r|_{W_A}),\galoisrc(\sigma^\natural)\bigr)$;
\item[(weights)] $\galoisHT_\tau=\{\frac{a-1}2-z:z\in\galoisinf_{\iota\tau}(\pi)\}$ and
  $\galoisHT_{\tau c}=\{\mathrm{wt}-h:h\in\galoisHT_\tau\}$ for every $\tau$;
\item[(crys)] for every place $w\mid\ell$ at which $\pi$ has a spherical local component, and every
  $A$ above $w$: $r|_{D_A}$ is crystalline, with $\tau$-Hodge--Tate weights $\galoisHT_\tau$ for every
  $\tau$ inducing $w$;
\item[(nice)] for every place $w\mid\ell$, every $A$ above $w$, every local component $\sigma$ of $\pi$
  at $w$ and every $\galoisrc$ with $\mathrm{Rc}(\galoisrc)$ such that $y=\galoisrc(\sigma^\natural)$ has
  $N_y=0$ and is either abelian, or tame and squarefree with $w$ $a$-large: $r|_{D_A}$ is de Rham and
  potentially semistable, with $\tau$-Hodge--Tate weights $\galoisHT_\tau$ for every $\tau$ inducing $w$,
  and $\sfname{AwayB}(\galoisWD(r|_{D_A}),y)$.
\end{description}
$\sfname{Row0b}(\mathrm{Rc})$ is the proposition: for every CM field $E$, algebraic closure $\Omega$,
prime $\ell$, isomorphism $\iota$, integer $a\ge1$, regular algebraic cuspidal automorphic
representation $\pi$ of $\GL_a(\A_E)$ and Hecke character $\chi$ of $E^+$ such that $(\pi,\chi)$ is
polarized, there are $r$, $\mathrm{wt}$ and $\galoisHT$ with
$\sfname{PropsB}_{\mathrm{Rc}}(\pi,\iota,r,\mathrm{wt},\galoisHT)$.
\end{definition}

\begin{theorem}[Galois representations with partial local--global compatibility]\Fmark\label{thm:row0b}
The proposition $\sfname{Row0b}(\galoisrecT)$ holds. That is, for every CM field $E$, prime $\ell$,
isomorphism $\iota$ and regular algebraic, polarized, cuspidal automorphic representation $(\pi,\chi)$ of
$\GL_a(\A_E)$, there are a continuous semisimple representation $r$ of $G_E$ on $\Qlbar^a$, an integer
$\mathrm{wt}$ and multisets $\galoisHT_\tau$ satisfying (away), (weights), (crys) and (nice) of
\cref{def:row0b} with respect to every local parameter map $\galoisrc$. In particular, if $w\nmid\ell$
is a finite place, $\sigma$ a local component of $\pi$ at $w$, $x=\galoisWD_\iota(r|_{W_w})$ and
$y=\galoisrc(\sigma^\natural)$, then $\dim x=a$; $x\cong y$ if $y$ is abelian with $N_y=0$; and $N_x=0$
and $x$, $y$ have the same characteristic polynomials on $I_w$ if $y$ is tame and squarefree with
$N_y=0$ and $E_w\cong\Q_g$ with $a=1$ or $g>16a^2$.
\formal{UpperTheory.Gal.Int.galoisLocalGlobalStrong_recParameter}
\end{theorem}
%P2PTR thm:row0b
\begin{proof}See \cref{sec:pf:geo}.
\end{proof}


At a place $w\nmid\ell$ where $\pi$ is unramified with Satake parameters $\alpha_1,\dots,\alpha_a$,
(L2) and (L3) make $y$ unramified, with geometric Frobenius eigenvalues $\Nm w^{(a-1)/2}\alpha_i$; so
$r$ is unramified at $w$ and $\iota\,\mathrm{charpoly}\,r(\Frob_w)=\prod_i(X-\Nm w^{(a-1)/2}\alpha_i)$
($\Frob_w$ geometric), which by the Chebotarev density and Brauer--Nesbitt theorems determines $r$. If
$\pi_\infty$ has the infinitesimal character of the trivial representation, then
$\galoisHT_\tau=\{0,1,\dots,a-1\}$, as in \cite{BarnetLambEtAl2014}. The proof gives in fact the variant
$\sfname{Row0}(\galoisrecT)$ without the condition ``$a=1$ or $g>16a^2$'' (\cref{thm:galois:reduction}).

\begin{littheorem}[{\cite[Thm.~2.1.1]{BarnetLambEtAl2014}}]\label{thm:galois:blggt}
Let $(\pi,\chi)$ be a regular algebraic, polarized, cuspidal automorphic representation of
$\GL_a(\A_E)$, $E$ a CM field (the theorem also covers totally real fields). There are a continuous
semisimple $r=r_{\ell,\iota}(\pi)\colon G_E\to\GL_a(\Qlbar)$ and $\mathrm{wt}\in\Z$ such that:
(1) $(r,\varepsilon_\ell^{1-a}r_{\ell,\iota}(\chi))$ is a totally odd polarized representation; in
particular $r^c\cong r^\vee\otimes\varepsilon_\ell^{1-a}r_{\ell,\iota}(\chi)|_{G_E}$;
(2) for every $w\nmid\ell$, $\iota\galoisWD(r|_{G_{E_w}})^{\mathrm{F\text{-}ss}}\cong
\mathrm{rec}_{E_w}(\pi_w\otimes|\det|^{(1-a)/2})$, and these representations are pure of weight
$\mathrm{wt}$;
(3) $r$ is de Rham at all places above $\ell$, with the Hodge--Tate weights determined by the weight
of $\pi_\infty$, and $\galoisHT_{\tau c}=\mathrm{wt}-\galoisHT_\tau$;
(4) if $w\mid\ell$ and $\pi_w$ has an Iwahori-fixed vector, then
$\iota\galoisWD(r|_{G_{E_w}})^{\mathrm{F\text{-}ss}}\cong\mathrm{rec}_{E_w}(\pi_w\otimes|\det|^{(1-a)/2})$;
in particular $r$ is semistable at $w$, and crystalline if $\pi_w$ is unramified.
\end{littheorem}

\begin{remark}[Comparison with \cref{thm:galois:blggt}]\label{rem:galois:blggt}
\Cref{thm:row0b} is weaker than \cref{thm:galois:blggt} in the following respects, and it is this
weaker statement that is used in the proof.
(a) Only CM fields are covered.
(b) There is no polarization of $r$ and no purity.
(c) At $w\nmid\ell$ the isomorphism $x\cong y$ is asserted only when $y$ is abelian with $N_y=0$, i.e.\
when $\sigma^\natural$ is an irreducible principal series (for instance when $\pi_w$ is unramified); then
$y=\mathrm{rec}_{E_w}(\sigma^\natural)$ by the compatibility of $\mathrm{rec}_{E_w}$ with class field theory
and the Langlands classification. For tame squarefree $y$ with $N_y=0$ at places with $E_w\cong\Q_g$,
$g>16a^2$ (or $a=1$), only $N_x=0$ and the characteristic polynomials on inertia are compared; in all
other cases (e.g.\ $N_y\ne0$, or $y$ wildly ramified and not abelian) only $\dim x=a$ is asserted.
(d) At $w\mid\ell$, crystallinity with the correct weights where $\pi_w$ is unramified agrees with
\cref{thm:galois:blggt}(4); elsewhere above $\ell$ the de Rham property and the comparison of
Weil--Deligne representations are asserted only for the parameters of (c), whereas
\cref{thm:galois:blggt} gives the de Rham property everywhere and the comparison whenever $\pi_w$ has an
Iwahori-fixed vector (including $N\ne0$).
(e) Compatibility is required for every local parameter map instead of $\mathrm{rec}_{E_w}$; by
\cref{rem:galois:suffices} the parameters for which the clauses have content are values at good
representations, on which all local parameter maps agree.
Thus \cref{thm:row0b} contains the existence of $r_{\ell,\iota}(\pi)$ with compatibility at all
unramified places, the Hodge--Tate weights, crystallinity at the unramified places above $\ell$, and
partial compatibility at ramified places of abelian or tame type. (The development contains a formal
statement of \cref{thm:galois:blggt}, which is not used.)
\end{remark}

% ---------------------------------------------------------------------------
\subsection{Architecture of the proof}\label{sec:galois:arch}

The proof of \cref{thm:row0b} passes through totally definite unitary groups. For $n\ge1$ let
$\galoisU_n(\A_{E^+})=\{g\in\GL_n(\A_E):{}^t\bar g\,g=1\}$ (the unitary group of the standard hermitian
form, definite at every infinite place); an automorphic representation $U$ of $\galoisU_n$ is an
irreducible space of automorphic forms on $\galoisU_n(E^+)\backslash\galoisU_n(\A_{E^+})$. At places of $E$
split over $E^+$ the local groups are $\GL_n(E_w)$, so $U$ has local components and Satake parameters;
at $\tau\colon E\to\C$ it has \emph{weights}, the dominant $\lambda\in\Z^n$ such that the algebraic
representation of highest weight $\lambda$ occurs in $U$ at $\tau$.

\begin{definition}[Weakly attached representations; $\sfname{CoreAll}$]\label{def:galois:core}
Let $(\pi,\chi)$ be as in $\sfname{Row0b}$. A representation $r\colon G_E\to\GL_a(\Qlbar)$ is
\emph{weakly attached} to $\pi$ if it is continuous and semisimple and, for almost all $w$, $\pi_w$ is
unramified, $r$ is unramified at $w$ and
$\iota\,\mathrm{charpoly}\,r(\Frob_w)=\prod_i(X-\Nm w^{(a-1)/2}\alpha_i)$, the $\alpha_i$ being the Satake
parameters of $\pi_w$. $\sfname{CoreAll}$ asserts: every such $\pi$ has a weakly attached $r$; and every
weakly attached $r$ satisfies, at every place $w$ where $\pi_w$ is unramified: if $w\nmid\ell$, the
equality just displayed; if $w\mid\ell$, $r|_{D_w}$ is crystalline with the Hodge--Tate weights of
(weights), and geometric Frobenius elements have characteristic polynomial
$\prod_i(X-\Nm w^{(a-1)/2}\alpha_i)$ on $\galoisWD(r|_{D_w})$.
\end{definition}

\begin{theorem}[Reduction to theta lifting and geometry]\Fmark\label{thm:galois:reduction}
\begin{enumerate}[label=(\alph*),nosep]
\item $\sfname{CoreAll}$ implies $\sfname{Row0}(\galoisrecT)$, the strengthening of
  $\sfname{Row0b}(\galoisrecT)$ in which ``$a$-large'' is replaced by ``$E_w\cong\Q_g$ for a prime $g$'';
  in particular it implies $\sfname{Row0b}(\galoisrecT)$.
\item $\sfname{GeoOne}\wedge\sfname{TLiftK}$ (\cref{def:galois:geoone,def:galois:tliftk}) implies
  $\sfname{CoreAll}$.
\end{enumerate}
Consequently \cref{thm:row0b} follows from \cref{thm:galois:tlift,thm:geoone}.
\formal{UpperTheory.Gal.Int.galoisLocalGlobal_of_geometry_lift}
\end{theorem}

\begin{proofoutline}
(a) Let $r$ be weakly attached to $\pi$. Comparing the infinitesimal characters of $\pi^c$ and
$\pi^\vee\otimes(\chi\circ\Nm\circ\det)$ through the invariant pairing gives $\mathrm{wt}\in\Z$ with
$\galoisinf_{\tau c}(\pi)=\{a-1-\mathrm{wt}-z:z\in\galoisinf_\tau(\pi)\}$, whence (weights). At unramified
places, (away) and (crys) are part of $\sfname{CoreAll}$. If $y$ (notation of \cref{thm:row0b}) is
abelian with $N_y=0$, then
$\sigma^\natural$ is an irreducible principal series (by (L1)--(L4) and the genericity of local
components of cuspidal representations \cite{Shalika1974}); for each $g\in W_w$ of nonzero degree,
solvable base change \cite{ArthurClozel1989} to a CM extension $M/E$ with $g\in G_M$ gives a cuspidal
representation of $\GL_a(\A_M)$, unramified at a place above $w$, to which $r|_{G_M}$ is weakly
attached, and $\sfname{CoreAll}$ for it computes the characteristic polynomial of $r(g)$; as elements of
nonzero degree determine characteristic polynomials on $W_w$, $x=y$. In the tame case $E_w\cong\Q_g$
and $\sigma$ has depth zero (\cref{rem:galois:suffices}); an unramified base change of suitable degree
turns the depth-zero supercuspidal representations in its support into principal series (the supports
being controlled through Shintani lifting for finite general linear groups, (L12) and
\cref{thm:rectexists}), the abelian case applies there, and characteristic polynomials on inertia do
not change under unramified base change. The same arguments with the clauses of $\sfname{CoreAll}$ at
$\ell$ give (nice).

(b) The proof has six steps.
\begin{enumerate}[label=(\roman*),nosep]
\item \emph{Seeds.} For odd $a$ and an automorphic representation $U$ of $\galoisU_a$ with a
  supercuspidal type at a place $w_0\mid\ell$, lacunary weight at a place $v'\mid\ell$ and generic
  Satake parameters at auxiliary places, $\sfname{GeoOne}$ provides representations whose Frobenius
  elements are annihilated by the Hecke polynomials of $U$. From them one extracts an $a$-dimensional
  continuous semisimple $r'$ whose Frobenius characteristic polynomials are the Hecke polynomials on
  an open dense set of Frobenius elements, using the lacunary Hodge--Tate weights, a genericity
  polynomial in the Satake parameters and determinant characters.
\item \emph{Congruences.} At a fixed level, the Hecke eigensystem of every automorphic representation
  of $\galoisU_a$ ($a$ odd, under field conditions) is a limit, for the topology of polynomial
  approximation, of eigensystems of seeds (density of generic Satake parameters, global Hecke
  operators of Plancherel type, $\ell$-adic families of lacunary weights).
\item \emph{Limits.} By the method of pseudo-characters \cite{Taylor1991,BellaicheChenevier2009,Chenevier2014}
  the seed representations converge to a continuous semisimple $a$-dimensional representation
  attached to $U$; crystallinity with bounded weights and the crystalline Frobenius survive in the
  limit (via Breuil--Kisin lattices, compare \cite{Berger2004,Kisin2008}).
\item \emph{From odd to arbitrary rank.} For even $n$, $\sfname{TLiftK}$ is applied twice, with
  $t=2(d_0+n)+1$ and two pairs of splitting characters, lifting $U$ to $U_1,U_2$ on $\galoisU_{n+t}$
  ($n+t$ odd). By (ii) of \cref{def:galois:tliftk} and the unramified type II correspondence, the Satake
  parameters of $U_k$ at split unramified places are those of $U$, twisted by a fixed algebraic Hecke
  character, together with the values of $t$ auxiliary algebraic Hecke characters, the two auxiliary
  families being disjoint. Removing the $\ell$-adic characters \cite{Serre1968} of the auxiliary
  characters from the representations of $U_1,U_2$ and twisting back gives the representation for $U$.
\item \emph{From unitary groups to $\GL_a$.} A descent of conjugate self-dual cuspidal representations
  of $\GL_a(\A_E)$ to $\galoisU_a$ matching local components at split places, in the style of
  \cite{Labesse2011,ClozelHarrisLabesse2011} (proved in the development from the Jacquet--Shalika bound
  and multiplicity one), transfers the representations to $\GL_a$.
\item \emph{Removing the conditions.} The conditions on the field and the restriction to split places
  are removed by base change to infinitely many quadratic CM extensions \cite{ArthurClozel1989}, with
  \cite{JacquetShalika1981,GelfandKazhdan1975}, and patching; polarized representations are reduced to
  conjugate self-dual ones by twisting with an algebraic Hecke character.
\end{enumerate}
Further published inputs are \cite{BernsteinZelevinsky1977,Zelevinsky1980,Tadic1986} and the theorem
that de Rham representations are potentially semistable \cite{Berger2002}; all of these are proved in
the formal development.

%P2PTR thm:galois:reduction
A complete proof is given in \cref{sec:pf:galred}.
\end{proofoutline}

% ---------------------------------------------------------------------------
\subsection{Theta lifts between definite unitary groups}\label{sec:galois:theta}

Let $\eta$ be the quadratic character of $\A_{E^+}^\times$ attached to $E/E^+$. A unitary Hecke
character $\chi$ of $E$ has \emph{archimedean exponents} $m=(m(w))_{w\mid\infty}$ with respect to a CM
type $(\tau_w)_{w\mid\infty}$ if $\chi(z)=\prod_w(\tau_w(z)/|\tau_w(z)|)^{m(w)}$ for archimedean $z$ near
$1$.

\begin{definition}[$\sfname{TLiftK}$]\label{def:galois:tliftk}
$\sfname{TLiftK}$ is the following statement. Let $E$ be a CM field such that $E/E^+$ is unramified at
all finite places, $E$ contains $\sqrt{-d}$ for some integer $d>0$, and every place of $E^+$ inert in
$E$ is prime to $2$ and has norm $>3$; fix a CM type $(\tau_w)$ and a standard additive character of
$\A_{E^+}/E^+$, unramified at the inert places. Let $n\ge1$ and let $X$ be the hermitian space $E^n$
with form $\sum_ix_i\bar y_i$ and a self-dual lattice. Then there is $d_0$ such that for every $d\ge
d_0$, every skew-hermitian space $W=(E^d,\delta\sum_ix_i\bar y_i)$ ($\bar\delta=-\delta$) definite at all
$\tau_w$ and with a self-dual lattice, all unitary Hecke characters $\chi_X,\chi_W$ of $E$ with
$\chi_X|_{\A_{E^+}^\times}=\eta^n$, $\chi_W|_{\A_{E^+}^\times}=\eta^d$ and archimedean exponents
$m_X,m_W$, and every automorphic representation $U$ of $\galoisU_n$ with weights $\lambda_w$ at $\tau_w$
such that $d+m_W(w)\le2\lambda_{w,k}$ for all $w$ and $k$, there is an automorphic representation
$U'$ of $\galoisU_d$ such that:
\begin{enumerate}[label=(\roman*),nosep]
\item for every $w\mid\infty$, $U'$ has a weight $\lambda'_w$ at $\tau_w$ with
  $2\lambda'_{w,k}=2\lambda_{w,k}-d-m_W(w)+n+m_X(w)$ for $k<n$ and $2\lambda'_{w,k}=n+m_X(w)$ for $n\le k<d$;
\item for every place $v$ of $E^+$ split in $E$ and all local components $U_v$, $U'_v$ of $U$, $U'$
  (representations of $\GL_n(E^+_v)$, $\GL_d(E^+_v)$), there is a nonzero map $\beta$ from the
  Schwartz space $S(M_{n\times d}(E^+_v))$ to the antilinear maps $U_v\to U'_v$ with
  $\beta(\omega_v(g,h)\Phi)\circ U_v(g)=U'_v(h)\circ\beta(\Phi)$, where $\omega_v$ is the action of
  $\GL_n\times\GL_d$ by translations twisted by $\chi_{W,v}^{-1}$, $\chi_{X,v}$, $|\det|^{-d/2}$ and
  $|\det|^{-n/2}$; that is, $(U_v,U'_v)$ occur in the type II local theta correspondence
  \cite{Minguez2008}.
\end{enumerate}
\end{definition}

\begin{theorem}[Theta lifts between definite unitary groups]\Fmark\label{thm:galois:tlift}
$\sfname{TLiftK}$ holds. \formal{UpperTheory.TP.K.tLiftK}
\end{theorem}
%P2PTR thm:galois:tlift
\begin{proof}See \cref{sec:pf:tlift}.
\end{proof}


Results of this type are classical: non-vanishing of theta lifts in the stable range via the Rallis
inner product \cite{Rallis1984,Li1992,HarrisKudlaSweet1996}, the correspondence for compact dual pairs
\cite{KashiwaraVergne1978,Howe1989}, and type II local theta correspondence \cite{Minguez2008}. The
statement in this form, with its threshold $d_0$, is specific to the development.

\begin{proofoutline}
Let $\omega$ be the Weil representation of $\galoisU(X)(\A)\times\galoisU(W)(\A)$ on the Schwartz space
of $M_{n\times d}(\A_{E^+})$, attached to the additive character and split by $\chi_X,\chi_W$
\cite{Weil1964,Kudla1994}. For $\Phi$ which is a polynomial Fock vector at the infinite places and
$\varphi\in U$, the theta lift
\[
  \theta(\Phi,\varphi)(h)=\int_{\galoisU(X)(E^+)\backslash\galoisU(X)(\A)}\ \sum_{\xi\in M_{n\times d}(E^+)}
  \omega(g,h)\Phi(\xi)\,\overline{\varphi(g)}\,dg
\]
is an automorphic form on $\galoisU(W)\cong\galoisU_d$, depending equivariantly on $(\Phi,\varphi)$.
(i) \emph{Archimedean types.} In the Fock model \cite{KashiwaraVergne1978,Howe1989} the irreducible
representation of highest weight $\lambda_w$ maps equivariantly into the Fock space, hitting a joint
highest weight vector whose type for $\galoisU(d)$ has highest weight $\lambda'_w$; here
$d+m_W\le2\lambda_w$ is used.
(ii) \emph{Non-vanishing.} Integrability of a height function on the doubled group
$G_2=\galoisU(X\oplus(-X))$ along $\galoisU(X)(\A)$ gives an exponent $t_0$, and one takes
$d_0=\max(2n,N)$ with $N>2t_0$. For $d\ge d_0$ the doubling integral
$\int_{\galoisU(X)(\A)}\langle\omega(g)\Phi,\Phi\rangle\overline{\langle\varphi(\cdot\,g),\varphi\rangle}\,dg$
converges absolutely and vanishes if $\theta(\Phi,\varphi)$ does (Rallis inner product); this follows from
the unfolding of the doubled theta integral and the Siegel--Weil formula on $G_2$ in the convergent
range $n<d$,
$E(x,\tfrac{d-n}2)=\vol^{-1}\int_{\galoisU(W)(E^+)\backslash\galoisU(W)(\A)}\theta(x,h;\Phi_1\otimes\bar\Phi_2)\,dh$
\cite{Weil1965,KudlaRallis1988,KudlaRallis1994,Ichino2004}. The doubling integral is nonzero for
suitable $\Phi$ (bump functions at finitely many places, a tail estimate, and Schur orthogonality at
the compact archimedean places).
(iii) \emph{Reading off the lift.} A nonzero lift has a nonzero equivariant projection to an
automorphic representation $U'$; $U'$ has the weights $\lambda'_w$ by (i), and at split places the
equivariance of the lift produces the map $\beta$. The Weil representation, the Siegel--Weil formula
in the form used and the Fock model are proved in the formal development.
\end{proofoutline}

% ---------------------------------------------------------------------------
\subsection{The geometric realization}\label{sec:galois:geo}

\begin{definition}[$\sfname{GeoOne}$]\label{def:galois:geoone}
Let $a$ be odd, $d\ge1$ and $x\in E$ with $x^2=-d$; assume that $-d$ is a square in $\Q_\ell$, that
$[E^+:\Q]$ is even, and that every place of $E$ not split over $E^+$ is prime to $\ell$ and unramified
over $E^+$. A place $w$ of $E$ is \emph{$d$-usable} if it is split over $E^+$ and $-d$ is a square in
$\Q_q$, $q$ the prime below $w$. Let $w_0,v'$ be places above $\ell$ lying over different places of
$E^+$ and over the same place of $\Q(x)$, let $S,T$ be finite sets of places, $T$ without split
places, and let $U$ be an automorphic representation of $\galoisU_a$ which is unramified at the
non-split places outside $T$; contains at $w_0$ a \emph{supercuspidal type of $\ell$-power order} (a
character $\chi\circ\vartheta$ of an open compact torsion-free $J\subset\GL_a(E_{w_0})$, with
$\vartheta\colon J\to\Z/\ell^m$ surjective and $\chi\ne1$, such that every smooth irreducible
representation containing it is supercuspidal); is spherical at $v'$ and at the split places lying
over places of $E^+$ different from those below $w_0$ and $S$; and has a \emph{lacunary weight} at
$v'$: for some $\tau\colon E\to\C$ inducing $v'$ via $\iota$ and some weight $\lambda$ of $U$ at $\tau$, no
two disjoint nonempty subsets of $\{k-\lambda_k:0\le k<a\}$ have the same average. For Satake
parameters $B$ of $U$ at $w$ put $P_{w,B}=\prod_{b\in B}(X-\Nm w^{(a-1)/2}b)$. Call a continuous
semisimple $\rho\colon G_E\to\GL_{a'}(\Qlbar)$ (of any dimension $a'$) \emph{basic} if (1) for all but
finitely many $d$-usable $w\nmid\ell$ at which $U$ is spherical with Satake parameters $B$, $\rho$ is
unramified at $w$ and $(\iota^{-1}P_{w,B})(\rho(\Frob_w))=0$; and (2) at some place above $\ell$, $\rho$ is
de Rham and, for some $\tau$, its $\tau$-Hodge--Tate multiset contains a lacunary set of $a$ integers.
$\sfname{GeoOne}$ asserts that under these hypotheses:
\begin{enumerate}[label=(G\arabic*),nosep]
\item a basic $\rho$ exists;
\item for every $d$-usable $w\nmid\ell$ at which $U$ is spherical, there is a basic $\rho$ which has
  the property of (1) at $w$;
\item for every $w\mid\ell$ at which $U$ is spherical, lying over the same place of $\Q(x)$ as $v'$ and
  over a place of $E^+$ different from those below $w_0$ and $v'$: there is a basic $\rho$ crystalline at
  $w$ whose $\tau$-Hodge--Tate weights contain $\{k-\lambda_k\}$ for every weight $\lambda$ of $U$ at
  every $\tau$ inducing $w$; and there is a basic $\rho$ crystalline at $w$ such that, for some element
  of $W_w$ of degree $-1$, its characteristic polynomial on $\galoisWD(\rho|_{D_w})$ divides a power of
  $P_{w,B}$.
\end{enumerate}
\end{definition}

\begin{theorem}[Geometric realization]\Fmark\label{thm:geoone}
$\sfname{GeoOne}$ holds. \formal{UpperTheory.Gal.Int.Geo.geoOne_holds}
\end{theorem}
%P2PTR thm:geoone
\begin{proof}See \cref{sec:pf:geo}.
\end{proof}


Results of this kind, for automorphic representations of definite unitary groups that are
supercuspidal at some place, are classical \cite{Kottwitz1990,Kottwitz1992,Clozel1991,HarrisTaylor2001};
the passage from the weak conclusions of $\sfname{GeoOne}$ to $a$-dimensional representations is step
(i) of the proof of \cref{thm:galois:reduction}(b). The least standard step of the proof uses the
following holomorphic Lefschetz formula for finite \'etale correspondences. For $c_1=\mathrm{id}$ it is the
holomorphic Lefschetz formula of Atiyah and Bott \cite{AtiyahBott1968,ToledoTong1978}. We do not know a published
proof for correspondences: a holomorphic Lefschetz formula for correspondences, going back to a question of
Shimura, is discussed in \cite{Tu2023} in the form of conjectures and an open problem, which includes the
computation of traces of Hecke correspondences. The development proves the formula below algebraically,
through \v{C}ech cohomology and a ``chart-lift'' theorem; the proof is given in \cref{sec:pf:geo}.

\begin{theorem}[Holomorphic Lefschetz formula for finite \'etale correspondences]\Fmark\label{thm:galois:hlf}
Let $Z$ be a smooth projective variety of dimension $\delta$ over $\C$ and $c_1,c_2\colon Y\to Z$ finite
\'etale morphisms. Assume that the set $\mathcal F$ of $y\in Y(\C)$ with $c_1(y)=c_2(y)$ is finite and
that for each $y\in\mathcal F$ the matrix $A_y\in M_\delta(\C)$ of the cotangent map of the local
self-map $c_2\circ c_1^{-1}$ at $y$ satisfies $\det(1-A_y)\ne0$. Then the spaces $H^b(Z,\Omega^p)$ are
finite-dimensional and, for every $p$,
\[
  \sum_{b=0}^{\delta}(-1)^b\tr\bigl((c_1)_*(c_2)^*\mid H^b(Z,\Omega^p)\bigr)
  =\sum_{y\in\mathcal F}\frac{\tr\bigl(\bigwedge^pA_y\bigr)}{\det(1-A_y)}.
\]
\formal{UpperTheory.Cb.lefschetz_fixedPoint_finite_etale_correspondence}
\end{theorem}

\begin{proofoutline}
The proof of \cref{thm:geoone} has three layers: group-theoretic \emph{plans} produce geometric
\emph{packages}, which produce the representations through \emph{step lemmas}.

\emph{(i) Packages and step lemmas.} A \emph{nice Hecke tower} over a field $k$ is a smooth projective
$k$-scheme $Z$ with correspondences $Z\xleftarrow{s_g}Y_g\xrightarrow{t_g}Z$ ($g\in\Gamma$), the $s_g$
finite \'etale; an eigensystem $\theta\colon\Gamma\to\Qlbar$ cuts out the $\theta$-eigenspace of the
operators $(s_g)_*(t_g)^*$ on the $\ell$-adic cohomology of $Z_\Omega$. A package for $U$ is a nice tower
over $E$ with an eigensystem, an algebraic Hecke character $\kappa$ of $E$, models over $\OO_w$ with
congruence data at the relevant places, and Hodge occurrences (see (iv)).
The representation is $\rho=\rho_0\otimes L(\kappa)$, with $\rho_0$ the semisimplified
$\theta$-eigenspace and $L(\kappa)$ the $\ell$-adic character of $\kappa$ \cite{Serre1968}. Its
properties come from step lemmas proved for arbitrary nice towers: the eigenspace is a
finite-dimensional continuous Galois representation; at $w\nmid\ell$, a model over $\OO_w$ on which
correspondences $T_0,\dots,T_n$ with eigenvalues $\vartheta_j$ satisfy the congruence relation of (iii)
makes $\rho_0$ unramified, with geometric Frobenius annihilated by the Hecke polynomial
$\sum_j(-1)^jq^{\binom j2}\vartheta_jX^{n-j}$ (specialization isomorphisms from proper and smooth base
change \cite{ArtinGrothendieckVerdier1972}); at $w\mid\ell$ with good reduction, $\rho_0$ is
crystalline (crystalline comparison theorem
\cite{FontaineMessing1987,Faltings1989,Tsuji1999,BhattMorrowScholze2018}), and with a congruence
relation the eigenvalues of the crystalline Frobenius are roots of the Hecke polynomial (compare
\cite{KatzMessing1974}); finally every $p$ for which $\theta$ occurs in $H^\bullet(Z_\tau,\Omega^p)$, $\tau$
inducing $w$, is a $\tau$-Hodge--Tate weight of $\rho_0$ (Hodge--Tate decomposition
\cite{Faltings1988,Scholze2013phodge}), which with the lacunary weight at $v'$ gives (2) of the
definition of basic representations.

\emph{(ii) Kuga--Sato towers over models of Harris--Taylor type.} For an auxiliary prime $p$ one uses a
PEL datum of Harris--Taylor type \cite{Kottwitz1992,HarrisTaylor2001}: $E$ with an imaginary quadratic
subfield in which $p$ splits, the lattice $\OO_E^a$ with an alternating form, and a unitary similitude
group $G$ of signature $(1,a-1)$ at one infinite place and definite at the others (unlike in
\cite{HarrisTaylor2001} the endomorphism algebra is $E$ itself, and compactness comes from the
definiteness at the other infinite places). The moduli of abelian schemes with $\OO_E$-action,
polarization, Kottwitz's determinant condition and level $N\ge3$ prime to $p$ are smooth
quasi-projective schemes $M_N$ of relative dimension $a-1$ \cite{Kottwitz1992}, proper and hence
projective here (compare \cite{Lan2013}); the development proves this using Mumford's
representability theorem, Zarhin's trick, Grothendieck--Messing theory and semistable reduction
\cite{MumfordFogartyKirwan1994,Zarhin1977,Messing1972,Grothendieck1972}. The towers consist of the
Kuga--Sato varieties $Z_N=\mathcal A_N^k$ (fibre powers of the universal abelian schemes) with the Hecke
correspondences of double cosets (compare \cite{Deligne1971,Scholl1990}); at places above primes split
in the imaginary quadratic field, integral models with additional \'etale level structure provide the
smooth projective models needed in (i).

\emph{(iii) A pointwise congruence relation.} Let $T_0,\dots,T_a$ be the correspondences of isogenies
of type $j$ at such a place $w$, $q=\Nm w$ and $\phi$ the $q$-Frobenius. For every point $y$ of the
integral model over a generic point of the special fibre, the multisets of reductions
\[
  \sum_{j\ \mathrm{even}}q^{\binom j2}\bigl\{\phi^{a-j}(\bar z):z\in T_j(y)\bigr\}
  \quad\text{and}\quad
  \sum_{j\ \mathrm{odd}}q^{\binom j2}\bigl\{\phi^{a-j}(\bar z):z\in T_j(y)\bigr\}
\]
coincide: a pointwise form of the Eichler--Shimura relation $\sum_j(-1)^jq^{\binom j2}T_j\Frob^{a-j}=0$
(compare \cite{Wedhorn2000}). Such points are ordinary at $w$ (compare \cite{Wedhorn1999}), and at an
ordinary point the $T_j$-images correspond to the $j$-dimensional subspaces of an $a$-dimensional
$\F_q$-space with a distinguished line, those containing the line being Frobenius twists; a
$q$-binomial count gives the identity.

\emph{(iv) Hodge occurrence.} The eigensystem of $U$ occurs in $H^\bullet(Z_N,\Omega^{h_i})$ for the
degrees $h_i$ determined by the weight of $U$ (after the twist by $L(\kappa)$ they become the
Hodge--Tate weights $\{k-\lambda_k\}$). The alternating traces on $H^b(Z_N,\Omega^p)$ of the Hecke
correspondences composed with fibrewise multiplications $[\alpha]$ are computed by
\cref{thm:galois:hlf}: the fixed points are CM points (regularity of the Hecke element forces complex
multiplication), with local terms given by the eigenvalues of the CM endomorphisms. The weighted count
of fixed points equals a sum over conjugacy classes in the definite inner forms of $G$, which equals
the trace of the Hecke operator on algebraic modular forms \cite{Gross1999} of these definite groups
(compare \cite{Kottwitz1992}; finiteness of class numbers \cite{Borel1963} is used). A Scholl-type combination $\sum_\alpha c_\alpha[\alpha]$
turns the result into $(\text{definite trace})\cdot\sum_iy^{h_i}$; the definite trace is nonzero for
suitable test functions because the eigensystem of $U$ occurs in algebraic modular forms of the
definite groups, and since the $h_i$ are distinct, comparing coefficients yields a nonzero Hecke
eigenvector in $H^\bullet(Z_N,\Omega^{h_i})$.

\emph{(v) Plans.} The automorphic input of the packages (eigensystem, test functions, Hecke data at
prescribed places, $\kappa$) is produced group-theoretically from $U$, by identifying the finite adelic
points of $G$ away from $p$ with those of the definite unitary similitude group of $E$; no comparison
of trace formulas between non-isomorphic groups is needed.

Most of the formal proof is foundational material proved in the development (\'etale cohomology,
$p$-adic comparison theorems, integral models, \cref{thm:galois:hlf}); the architecture, the
congruence relation on Kuga--Sato varieties, the fixed-point computation and the proof of Hodge
occurrence are, as far as we know, new.

%P2PTR thm:galois:hlf
A complete proof is given in \cref{sec:pf:geo}.
\end{proofoutline}
